\renewcommand{\ElicitationDate}{2026.11.19}

\section{Elicitation \ElicitationDate}

\begin{tabular}{ll}
\textbf{Date:} & \ElicitationDate \\
\textbf{Consultant:} & Hector Hill / Gy'ax \\
\textbf{Elicitors:} & Zahra, Andrew, Matthias \\
\textbf{Location:} & TFS elicitation room \\
\textbf{Context:} & Eighth elicitation for LING 431/531 Field Methods \\
\textbf{Audio:} & \texttt{November\_19\_2026\_HH.WAV} \\
\end{tabular}

\subsection{Matthias [00:00-37:25]}

[2:00]:

\textbf{Context: }You are invitied to a soccer game and there is a requirement that you bring at least three more people. In fact, you bring five friends.

HH: "When you're adding the \textit{g̱abii} you are putting 'three and three'."

\begin{exe}
    \ex \begin{xlist}
        \item  \glll da'witxwinhl gwilunhl ensiip'insxwaa\\
            da-'witxw-n=hl gwilun=hl en-siip'-in-sxw=aa\\
            \Com{}-come-\Ssg{}.\Ii{}=\Cn{} three.\Hum{}=\Cn{} \Nmlz{}-like-\Caus2{}-\Antip{}=\Q{} \\
            \glt "Did you bring three friends?" \hfill HH.\ElicitationDate{} (SF)
        \item \glll ee, da'witxw'yhl gwilunhl ensiip'insxw, iin ḵ'ap da'witxwhl xwisdinsulhl ensiip'insxw \\
            ee, da-'witxw-'y=hl gwilun=hl en-siip'-in-sxw, ii=n ḵ'ap da-'witxw=hl xwisdinsul=hl en-siip'-in-sxw \\
            yes, \Com{}-come-\Fsg{}.\Ii{}=\Cn{} three.\Hum{}=\Cn{} \Nmlz{}-like-\Caus2{}-\Antip{}=\Q{}, \Ccnj{}=\Fsg{}.\I{} \Verum{} \Com{}-come-\Fsg{}.\Ii{}=\Cn{} five.\Hum{} \Nmlz{}-like-\Caus2{}-\Antip{}=\Q{} \\
            \glt "yes, I brought three friends, and in fact I brought five friends." \hfill HH.\ElicitationDate{} (SF)
        \item \glll \# da'witxw'yhl gwilunhl g̱abiihl ensiip'insxw, iin ḵ'ap da'witxwhl xwisdinsulhl ensiip'insxw \\
            {} da-'witxw-'y=hl gwilun=hl g̱abii=hl en-siip'-in-sxw, ii=n ḵ'ap da-'witxw=hl xwisdinsul=hl en-siip'-in-sxw \\
            {} \Com{}-come-\Fsg{}.\Ii{}=\Cn{} three.\Hum{}=\Cn{} \Cnt{}=\Cn{} \Nmlz{}-like-\Caus2{}-\Antip{}=\Q{}, \Ccnj{}=\Fsg{}.\I{} \Verum{} \Com{}-come-\Fsg{}.\Ii{}=\Cn{} five.\Hum{} \Nmlz{}-like-\Caus2{}-\Antip{}=\Q{}\\
            \glt "I brought three friends (total), in fact I brought five friends." \hfill HH.\ElicitationDate{} (SF) \\
            Comment: "\textit{g̱abiihl} may be in the wrong place."
    \end{xlist}
\end{exe}

[4:01]: I think \textit{suwii} precedes \textit{k'wootxw} here.

\begin{exe}
    \ex \glll suwii k'wootxwinaa \\
    suwii k'woot-xw-n=aa \\
    away lost-\Val{}-\Ssg{}.\Ii{}=\Q{} \\
    \glt "Were you lost?" \hfill HH.\ElicitationDate{} (VF)
\end{exe}

[4:15]: HH: "it's ok to use \textit{g̱abii} when you're telling the amount".

[4:25]: This continuation wasn't totally accepted, but not rejected either. Preferred continuation at [5:20].

\begin{exe}
    \ex \glll ee, da'witxw'yhl gwilunhl ensiip'insxw, ? iin ḵ'ap da'witxwhl xwisdinsulhl g̱abiihl ensiip'insxw \\
            ee, da-'witxw-'y=hl gwilun=hl en-siip'-in-sxw, {} ii=n ḵ'ap da-'witxw=hl xwisdinsul=hl g̱abii=hl en-siip'-in-sxw \\
            yes, \Com{}-come-\Fsg{}.\Ii{}=\Cn{} three.\Hum{}=\Cn{} \Nmlz{}-like-\Caus2{}-\Antip{}=\Q{}, {} \Ccnj{}=\Fsg{}.\I{} \Verum{} \Com{}-come-\Fsg{}.\Ii{}=\Cn{} five.\Hum{} \Cnt{}=\Cn{} \Nmlz{}-like-\Caus2{}-\Antip{}=\Q{} \\
            \glt "Yes, I brougght three friends, in fact I brought five in total." \hfill HH.\ElicitationDate{} (SF)
\end{exe}

[5:20]:

\begin{exe}
    \ex \glll
    ... ii sag̱ayt xwsdins g̱abiihl da'witxw'y\\
    ... ii sag̱ayt xwsdins g̱abii=hl da-'witxw-'y\\
    ... \Ccnj{} together five \Cnt{}=\Cn{} \Com{}-come-\Fsg{}.\Ii{} \\
    \glt "... and together I brought five" Lit: "and together the amount that I brought is five." \hfill HH.\ElicitationDate{} (VF)
\end{exe}

[6:20]: \textbf{Conext: }You are invited to a soccer game, here there isa requirement that you bring exactly three people so that there are just enough people to play the game.

\begin{exe}
    \ex \begin{xlist}
        \item \glll da'witxwinhl gwilunhl ensiip'insxwaa\\
            da-'witxw-n=hl gwilun=hl en-siip'-in-sxw=aa\\
            \Com{}-come-\Ssg{}.\Ii{}=\Cn{} three.\Hum{}=\Cn{} \Nmlz{}-like-\Caus2{}-\Antip{}=\Q{} \\
            \glt "Did you bring three friends?" \hfill HH.\ElicitationDate{} (SF)
        \item \glll ee, (ḵ'ap) da'witxw'yhl gwilunhl ensiip'insxw \\
               ee, (ḵ'ap) da-'witxw-'y=hl gwilun=hl en-siip'-in-sxw \\
               yes, (\Verum{}) \Com{}-come-\Fsg{}.\Ii{}=\Cn{} three.\Hum{}=\Cn{} \Nmlz{}-like-\Caus2{}-\Antip{} \\
               \glt "Yes, I brought three friends." \hfill HH.\ElicitationDate{} (SF)
               Comment: \textit{ḵ'ap} is what you asked for, it's the total.
        \item \glll \# ee, ḵ'ap da'witxw'yhl gwilunhl g̱abiihl ensiip'insxw \\
               {} ee, ḵ'ap da-'witxw-'y=hl gwilun=hl g̱abii=hl en-siip'-in-sxw \\
               {} yes, \Verum{} \Com{}-come-\Fsg{}.\Ii{}=\Cn{} three.\Hum{}=\Cn{} \Cnt{}=\Cn{} \Nmlz{}-like-\Caus2{}-\Antip{} \\
               \glt Intended: "Yes, I brought three friends total." \hfill HH.\ElicitationDate{} (SF)
               Comment: not \textit{g̱abiihl}.
    \end{xlist}
\end{exe}

[7:25]: HH: "If you wanted to answer what he asked for:"

\begin{exe}
    \ex \begin{xlist}
            \item \glll
            ee, ḵ'ap dabekw'y gwilunhl ensiipinsxw\\
            ee, ḵ'ap da-bekw'-y gwilen=hl en-siip'-in-sxw\\
            yes, \Verum{} \Com{}-arrive-\Fsg{}.\Ii{} three.\Hum{}=\Cn{} \Nmlz{}-like-\Caus2{}-\Antip{} \\
            \glt "Yes, I actually brought three friends." \hfill HH.\ElicitationDate{} (VF)
            \item \glll
            ee, ḵ'ap dabekw'y gwilenhl ensiip'insxw hasag̱an \\
            ee, ḵ'ap da-bekw'-y gwilen=hl en-siip'-in-sxw hasaḵ-an\\
            yes, \Verum{} \Com{}-arrive-\Fsg{}.\Ii{} three.\Hum{}=\Cn{} \Nmlz{}-like-\Caus2{}-\Antip{} desire-\Ssg{}.\Ii{} \\
            \glt "Yes, I actually brought the three friends that you wanted." \hfill HH.\ElicitationDate{} (VF)
    \end{xlist}
\end{exe}

[9:30]: Can also simply answer:

\begin{exe}
    \ex \glll
    ee, gwilunhl dabegw'y\\
    ee, gwilun=hl da-bekw-'y\\
    yes, three.\Hum{}=\Cn{} \Com{}-arrive-\Fsg{}.\Ii{} \\
    \glt "Yes, I brought three." \hfill HH.\ElicitationDate{} (VF)
\end{exe}

[9:48]: HH unprompted, "the other one where you brought three and two: "

I'm not actually sure what this sentence means exactly. Is the first clause negated? Wouldn't we see \textit{nee dii} in that case?

\begin{exe}
    \ex \glll
    nee, dabegw'y gwilent, iin ḵ'ap dabekwhl bag̱adilt \\
    nee, da-begw-'y gwilen-t, ii=n ḵ'ap da-bekw=hl bag̱adil-t  \\
    no, \Com{}-arrive-\Fsg{}.\Ii{} three.\Hum{}-3, \Ccnj{}=\Fsg{}.\I{} \Com{}-arrive=\Cn{} two.\Hum{}-3 \\
    \glt "No, I brought three, and actually I brought two." \hfill HH.\ElicitationDate{} (VF)
\end{exe}

[10:05]: \textbf{Context: }In this case you can bring a maximimum of three people.

HH: Can't say \textit{ee} if you didn't bring the three.

\begin{exe}
    \ex \glll
    \# ee, da'witxw'yhl gwilunhl ensiip'insxw, iin ḵ'ap da'witxwhl ḵ'am bag̱adilhl ensiip'insxw  \\
    {} ee, da-'witxw-'y=hl gwilun=hl en-siip'-in-sxw, ii=n ḵ'ap da-'witxw=hl ḵ'am bag̱adil=hl en-siip'-in-sxw\\
    {} yes, \Com{}-come-\Fsg{}.\Ii{}=\Cn{} three.\Hum{}=\Cn{} \Nmlz{}-like-\Caus2{}-\Antip{}=\Q{}, \Ccnj{}=\Fsg{}.\I{} \Verum{} \Com{}-come-\Fsg{}.\Ii{}=\Cn{} only two.\Hum{} \Nmlz{}-like-\Caus2{}-\Antip{} \\
    \glt "Yes, I'm bringing friends, and actually I'm bringing only two." \hfill HH.\ElicitationDate{} (SF)
\end{exe}

[11:18]:

\begin{exe}
    \ex \glll
    nee, ḵ'am bag̱adilhl dabegw'y \\
    nee, ḵ'am bag̱adil=hl da-bekw-'y \\
    no, only two.\Hum{}=\Cn{} \Com{}-arrive-\Fsg{}.\Ii{} \\
    \glt "No, I only brought two." \hfill HH.\ElicitationDate{} (VF)
\end{exe}

[11:25]: I am unclear what exactly follows \textit{hli ky'ults} and \textit{aliisxw}. Probably something meaning "come", but I can't make it out.

\begin{exe}
    \ex \glll
    dim 'witxwhl hli ky'ult ? aliisxw\\
    dim 'witxw=hl hli ky'ul-it ? aliisxw \\
    \Prosp{} arrive=\Cn{} \Part{} one.\Hum{}-\Sx{} ? later \\
    \glt "The other one will come later." \hfill HH.\ElicitationDate{} (VF)
\end{exe}

[11:55]: Ideas for "maximum": \textit{hogax̱hl g̱abiit}, \textit{hogax̱hl g̱abiit}, \textit{sida}.

[12:10]: This could be used to ask for a maximum. In a way, we can then analyze \textit{ḵ'am} as closing off the numeral scale. I clearly hear a voiced \textit{gw} and also first person marking, but doesn't make sense unless the translation is as stated.

Maybe \textit{g̱abiihl} is dispreferred here because we already have \textit{ḵ'am}?

\begin{exe}
    \ex \glll
    ḵ'am gwile'lhl (*g̱abiihl) dim dabegw'y \\
    ḵ'am gwile'l=hl (*g̱abii=hl) dim da-begw-'y\\
    only three=\Cn{} (\Cnt{}=\Cn{}) \Prosp{} \Com{}-arrive-\Fsg{}.\Ii{} \\
    \glt "I will only bring three" \hfill HH.\ElicitationDate{} (VF)
\end{exe}

[13:07]:

\begin{exe}
    \ex \glll
    ḵ'am gwile'lhl ensiip'insxw dabegw'y \\
    ḵ'am gwile'l=hl en-siip'-in-sxw da-begw-'y \\
    only three=\Cn{} \Nmlz{}-like-\Caus2{}-\Antip{} \Com{}-arrive-\Fsg{}.\Ii{} \\
    \glt "I only brought three friends." \hfill HH.\ElicitationDate{} (VF)
\end{exe}

[13:12]: even after HH repeated it multiple times ([13:20]), I don't know what is after ky'ult.

I'm also kind of confused about the order of \textit{dim} and \textit{ii}... but this order is attested in the storybook.

\begin{exe}
    \ex \glll
    nee ḵ'am gilbil ensiip'insxw dabegw'y, dim ii kw'itxwhl hla ky'ult  ? aliisxw \\
    nee ḵ'am gilbil en-siip'-in-sxw da-begw-'y, dim ii kw'itxw[-t]=hl hla ky'ul-t ? aliisxw\\
    no only two  \Nmlz{}-like-\Caus2{}-\Antip{} \Com{}-arrive-\Fsg{}.\Ii{}, \Prosp{} \Ccnj{} come[-3]=\Cn{} \Part{} one.\Hum{}-3 ? later\\
    \glt "No I only brought two, the other one will come later." \hfill HH.\ElicitationDate{} (VF)
\end{exe}

[15:18]:

\begin{exe}
    \ex \glll
    da'aḵhlxwin dimim dabekwhl gwile'lhl ensiip'insxw \\
    da'aḵhlxw-i-n dim=m da-bekw[-t]=hl gwile'l=hl en-siip'-in-sxw\\
    able-\Tr{}-\Ssg{}.\Ii{} \Prosp{}=\Ssg{}.\I{} \Com{}-arrive[-3]=\Cn{} three=\Cn{} \Nmlz{}-like-\Caus2{}-\Antip{} \\
    \glt "You can bring three friends." \hfill HH.\ElicitationDate{} (VF)
\end{exe}

[15:35]: Is this a case of long-distance movement?? At least as per the literal translation that I am assuming. The morphology sure looks like \textsc{ox} for both clauses.

\begin{exe}
    \ex \glll
    ḵ'am gwile'lhl ensiip'insxw dim da'aḵhlxwin dim dabegwin \\
    ḵ'am gwile'l=hl en-siip'-in-sxw dim da'aḵhlxw-i-n dim da-bekw-i-n\\
    only three=\Cn{} \Nmlz{}-like-\Caus2{}-\Antip{} \Prosp{} able-\Tr{}-\Ssg{}.\Ii{} \Prosp{} \Com{}-arrive-\Tr{}-\Ssg{}.\Ii{}  \\
    \glt Lit: "only three are the friends that you will be able to bring." \hfill HH.\ElicitationDate{} (VF)
\end{exe}

[16:03]: Can add \textit{ii ḵ'ap} in here somewhere? Probably at the start to be "very precise".

\begin{exe}
    \ex \glll
    skidim ḵ'am gwile'lhl (ensiip'insxw) dim dabegwin \\
    skidim ḵ'am gwile'l=hl (en-siip'-in-sxw) dim da-begw-i-n\\
    should only three=\Cn{} (\Nmlz{}-like-\Caus2{}-\Antip{}) \Prosp{} \Com{}-arrive-\Tr{}-\Ssg{}.\Ii{} \\
    \glt Lit: "the friends that you will bring should only be three." \hfill HH.\ElicitationDate{} (VF)
\end{exe}

[16:50]: 

\begin{exe}
    \ex \glll
    hasag̱a'mhl gwilel ensiip'insxw, ii ḵ'ap da'aḵhlxwin dimin dabekwhl hlag̱ats'uulhl gat\\
    hasaḵ-'m=hl gwilel en-siip'-in-sxw, ii ḵ'ap da'aḵhlxw-i-n dim=m da-bekw[-t]=hl hla-g̱a-ts'uu=hl gat \\
    desire-\Fpl{}.\Ii{}=\Cn{} three \Nmlz{}-like-\Caus2{}-\Antip{}, \Ccnj{} \Verum{} able-\Tr{}-\Ssg{}.\Ii{} \Prosp{}=\Ssg{}.\I{} \Com-arrive[-3]=\Cn{} \Part{}-\Distr{}-other=\Cn{} person \\
    \glt "We only need three friends, but you can bring as much as you want (=other people)." \hfill HH.\ElicitationDate{} (VF)
\end{exe}

[17:30]: \textbf{Context: }I went on a hike, and when I was on the hike I heard from other hikers that they saw bears around. I am telling this to a friend.

HH: If we use \textit{hilt=hl} here, this means "now they are looking at the bear, \textit{hilt} [without \textit{=hl}] is past... \textit{=hl} is like the present." What to make of this

\begin{exe}
    \ex \glll
    hilt an ga'ahl smex \\
    hilt an[-t] ga'a=hl smex\\
    many \Ax[-3] see=\Cn{} bear\\
    \glt "Many saw a bear (=the ones who saw a bear are many)" \hfill HH.\ElicitationDate{} (SF)
\end{exe}

[18:55]: Here again, "if you're talking on the phone, you can say \textit{hilthl gat an ga'ahl smex}".

\begin{exe}
    \ex \glll
    hilt gat an ga'ahl smex \\
    hilt gat an[-t] ga'a=hl smex\\
    many person \Ax[-3] see=\Cn{} bear\\
    \glt "Many people saw a bear" \hfill HH.\ElicitationDate{} (VF)
\end{exe}


[19:38]: I hear both \textit{-in} and \textit{dip} here.

\begin{exe}
    \ex \glll
    yukw dimin dip ga'ahl smex\\
    yukw dim=n dip ga'a[-t]=hl smex \\
    \Prog{} \Prosp{}=\Fpl{}.\I{} \Fpl{}.\I{} see[-3]=\Cn{} bear\\
    \glt "We are going to see a bear." \hfill HH.\ElicitationDate{} (VF)
\end{exe}

[20:10]:

\begin{exe}
    \ex \glll
    hilt dim an ga'ahl smex\\
    hilt dim an[-t] ga'a[-t]=hl smex\\
    many \Prosp{} \Ax{}[-3]=\Cn{} see[-3]=\Cn{} bear\\
    \glt "Many will see a bear (=many are those who will see a bear)." \hfill HH.\ElicitationDate{} (VF)
\end{exe}

[21:35]: \textbf{Context: }You are a tour guide leading a group, you want to tell the group what we will see.

\begin{exe}
    \ex \glll
    hilt smex dim ga'a'm\\
    hilt smex dim ga'-a-'m\\
    many bear \Prosp{} see-\Tr{}-\Fpl{}.\Ii{}\\
    \glt "We will see many bears (many are the bears that we will see.)" \hfill HH.\ElicitationDate{} (VF)
\end{exe}

[22:20]: I am trying to get \Ax{} to work here, but Hector repeats it with \Ox{} as above.

Hector repeats this sentence, but I don't think he ever confirmed them as being possible, since he asks "is this what you want?".

\begin{exe}
    \ex \glll
    ? hilthl gat dim an ga'ahl smex\\
    {} hilt=hl fat dim an[-t] ga'a=hl smex\\
    {} many=\Cn{} person \Prosp{} \Ax{}[-3] see=\Cn{} bear\\
    \glt Intended: "Many people will see a bear." \hfill HH.\ElicitationDate{} (SF)
\end{exe}

[24:00]: some good phrases around this time about "can I saw this" (useful for eliciting in Gitksan), return to these.

[25:05]: \textbf{Context: }You went on a hike and you heard from two hikers that they saw bears.

I ask at [26:00] whether the presence of \textit{=hl} affects the tense at all for Hector. 

\begin{exe}
    \ex \glll
    bag̱adilhl gat an ga'ahl smex \\
    bag̱adilhl gat an[-t] ga'a=hl smex\\
    two.\Hum{}=\Cn{} person \Ax{}[-3] see=\Cn{} baer\\
    \glt "Two people saw a bear." \hfill HH.\ElicitationDate{} (SF)
\end{exe}\

[26:50]: Now HH says you can't use \textit{hilt gat}, even though in (488) he drops the \textit{=hl} as far as I can tell. 

[27:30]:

\begin{exe}
    \ex \glll
    hilthl gat begwit \\
    hilt=hl gat begw-it\\
    many=\Cn{} person arrive-\Sx{}\\
    \glt "A lot of people came." \hfill HH.\ElicitationDate{} (VF)
\end{exe}

[27:50]:

\begin{exe}
    \ex \glll
    hilthl gat dim begwit \\
    hilt=hl gat dim begw-it\\
    many=\Cn{} person \Prosp{} arrive-\Sx{}\\
    \glt "A lot of people are coming." \hfill HH.\ElicitationDate{} (VF)
\end{exe}

[27:58]: Once again HH repeats (488) without \textit{=hl}, so you can have \textit{hilt gat}...... 

Here HH confirms again his interpretation of present with \textit{=hl}.

[28:45]: This by itself isn't really accepted by Hector, "there's that one word that's tipping it to the side". 

\begin{exe}
    \ex \glll
    ? ga'ahl bag̱adilhl gathl smex\\
    {} ga'a=hl bag̱adil=hl gat=hl smex \\
    {} see=\Cn{} two.\Hum{}=\Cn{} person=\Cn{} bear\\
    \glt Intended: "Two people saw a bear." \hfill HH.\ElicitationDate{} (SF)
\end{exe}

[32:20]: I think this is accepted, but at the minimum it is unerstood.

\begin{exe}
    \ex \glll
    hindehl g̱abiihl gat an ga'ahl smex \\
    hinde=hl g̱abii=hl gat an[-t] ga'a=hl smex\\
    \Wh{}=\Cn{} \Cnt{}=\Cn{} person \Ax{}[-3] see=\Cn{} bear\\
    \glt "How many people saw bears?" \hfill HH.\ElicitationDate{} (SF)
\end{exe}

[32:40]: As answer to the above.

\begin{exe}
    \ex \glll
    ḵ'am bag̱adilhl gat an ga'ahl smex \\
    ḵ'am bag̱adil=hl gat an[-t] ga'a=hl smex\\
    only two.\Hum{}=\Cn{} person \Ax{}[-3] see=\Cn{} bear\\
    \glt "Only two people saw a bear" \hfill HH.\ElicitationDate{} (VF)
\end{exe}

[33:35]: I try to ask the question in English, "did the two people see the bears or hear the bears". Here we could answer \textit{ḵ'am ga'a} `only saw'. 

Otherwise, for a full answer, we answer exactly as above. Probably didn't do a good job focusing the see vs. hear.

[33:55]: Don't actually see \textit{x̱a'ni} in FV, but it is documented in the storybook.

Also, why is \textsc{3.i} double marked in the latter part?

\begin{exe}
    \ex \glll
    ḵ'am bag̱adilhl gat an ga'ahl smex, iit nee diit x̱a'nidiit 'nidiit \\
    ḵ'am bag̱adil=hl gat an[-t] ga'a=hl smex, ii=t nee dii=t x̱a'ni-diit 'nidiit \\
    only two.\Hum{}=\Cn{} person \Ax{}[-3.\Ii{}] see=\Cn{} bear, \Ccnj{}=3.\I{} \Neg{} \Foc=3.\I{} hear-\Tpl{}.\I{} \Tpl{}.\Iii{}\\
    \glt "Two people saw the bears but they didn't hear them." \hfill HH.\ElicitationDate{} (VF)
\end{exe}